\chapter{Wstęp}
\label{ch:wstep}

% W tym rozdziale cytujemy WYŁĄCZNIE źródła faktycznie posiadane i czytane:
%   - tatari2025splitcost  (IEEE JESTPE, 2025) — Artykuł 1, split-cost FS-MPC
%   - tatari2024mfbb       (Power Electronics and Drives) — Artykuł 2 BAZOWY, MF-BB
% Pozostałe twierdzenia merytoryczne mają oznaczenie "% TODO cytat: ..."
% z opisem czego brakuje. Bibliografię uzupełniamy systematycznie w miarę
% pozyskiwania źródeł (skany, IEEE Xplore, biblioteka PW).

\section{Wprowadzenie do tematyki}
\label{sec:wprowadzenie}

Dzisiejszy rozwój energetyki rozproszonej oraz dynamiczna ekspansja odnawialnych źródeł energii (OZE) powodują, że tradycyjny model scentralizowanej dystrybucji energii ustępuje miejsca koncepcji \emph{mikrosieci} -- lokalnie zarządzanych układów elektroenergetycznych łączących wytwarzanie, magazynowanie i~zużycie energii. Bardzo istotnym wariantem są mikrosieci prądu stałego (DC), które naturalnie odpowiadają charakterystyce pracy ogniw fotowoltaicznych, akumulatorów elektrochemicznych oraz nowoczesnych odbiorników półprzewodnikowych. Wykluczają one potrzebę wielokrotnego przetwarzania energii pomiędzy domenami AC i~DC, co bezpośrednio przekłada się na wyższą sprawność oraz prostszą architekturę regulacji.

Bardzo ważnym elementem każdej mikrosieci DC są \emph{dwukierunkowe przekształtniki energoelektroniczne}, które odpowiadają za kontrolowany przepływ energii pomiędzy szyną zbiorczą a~poszczególnymi modułami systemu -- źródłami odnawialnymi, magazynami energii czy odbiorami~\cite{erickson2020}. Jednym z najczęściej stosowanych rozwiązań w energoelektronice jest układ podwyższający (Boost) pracujący w topologii dwukierunkowej, który umożliwia elastyczną pracę w~szerokim zakresie napięć, oferując przy tym pełną swobodę w zakresie kierunku transferu mocy.

Zapewnienie stabilnego i dokładnego sterowania takimi systemami jest skomplikowanym zadaniem inżynierskim, ponieważ ich dynamika cechuje się silną nieliniowością oraz, w~niektórych punktach pracy, zjawiskiem fazy nieminimalnej (ang. \emph{non-minimum phase}, NMP), polegającym na początkowej, odwrotnej do zamierzonej, reakcji napięcia wyjściowego na zmiany sygnału sterującego~\cite{tatari2024mfbb}.

\section{Motywacja podjęcia tematu}
\label{sec:motywacja}

Obecne podejście do projektowania w~obszarze energoelektroniki opiera się w~znacznej mierze na wykorzystaniu komercyjnych, zamkniętych pakietów symulacyjnych, takich jak PSIM czy MATLAB/Simulink. Narzędzia te są bardzo cenione za wysoką precyzję oraz bogate biblioteki komponentów, jednak ich komercyjny i~zamknięty charakter stwarza istotne ograniczenia z punktu widzenia współczesnych standardów inżynierii oprogramowania.

Po pierwsze, ich monolityczna struktura ogranicza swobodę modyfikacji algorytmów numerycznych, a także utrudnia bezpośrednią integrację ze współczesnymi bibliotekami obliczeniowymi oraz algorytmami sztucznej inteligencji. Po drugie, wysokie koszty licencji tworzą barierę utrudniającą reprodukcję wyników przez niezależne zespoły badawcze. Wreszcie, projekty tworzone w~zamkniętych pakietach symulacyjnych z~trudem poddają się standardowym praktykom inżynierii oprogramowania, takim jak wersjonowanie kodu, automatyzacja testów jednostkowych czy ciągła integracja (CI/CD).

Alternatywę stanowi język Python, który wraz z dojrzałym ekosystemem bibliotek numerycznych (NumPy~\cite{harris2020numpy}, Matplotlib~\cite{hunter2007matplotlib}) oraz optymalizacyjnych (PySwarms~\cite{miranda2018pyswarms}) tworzy otwarte, w~pełni rozszerzalne środowisko badawcze.

Przeniesienie modeli energoelektronicznych do takiego ekosystemu otwiera możliwość bezpośredniej integracji z~metodami sztucznej inteligencji oraz pełnej automatyzacji procesu strojenia sterowników. Niniejsza praca podejmuje to wyzwanie poprzez opracowanie modelu, który nie tylko jest poprawny fizycznie, lecz również spełnia wysokie standardy inżynierii oprogramowania, takie jak modularność, testowalność i~rozszerzalność powstałego kodu źródłowego.

\section{Cel pracy}
\label{sec:cel}

\textbf{Celem pracy jest} zaprojektowanie oraz zaimplementowanie autorskiego \emph{Bliźniaka Cyfrowego} (ang. \emph{Digital Twin}) dwukierunkowego przekształtnika prądu stałego typu Boost w~języku Python, zintegrowanego z~bezmodelowym algorytmem sterowania typu Bang-Bang oraz metaheurystycznym algorytmem strojenia opartym o~Optymalizację Rojem Cząstek (ang. \emph{Particle Swarm Optimization}, PSO).

Tak sformułowany cel obejmuje trzy dopełniające się aspekty:

\begin{itemize}
	\item \textbf{aspekt fizyczny} -- opracowanie modelu wysokiej wierności, który precyzyjnie oddaje dynamikę przekształtnika (w tym efekt NMP). Wiarygodność modelu została zweryfikowana poprzez analizę porównawczą z wynikami uzyskanymi w profesjonalnym środowisku PSIM;
	\item \textbf{aspekt programistyczny} -- zaprojektowanie obiektowej, modularnej architektury kodu zgodnej z~zasadami inżynierii oprogramowania, gwarantującej rozszerzalność i~reprodukowalność eksperymentów;
	\item \textbf{aspekt badawczy} -- zastosowanie algorytmu PSO do automatycznego doboru parametrów sterownika Bang-Bang oraz wykazanie przewagi tak strojonego układu nad strojeniem ręcznym, w~szczególności w~warunkach ograniczonej pojemności filtra wyjściowego.
\end{itemize}

\section{Zakres pracy}
\label{sec:zakres}

W~ramach realizacji postawionego celu zaplanowano i~wykonano następujące zadania programistyczne oraz badawcze:

\begin{enumerate}
	\item Inżynieria odwrotna referencyjnego obwodu Boost zaimplementowanego w~oprogramowaniu PSIM, prowadząca do sformułowania i wyprowadzenia kompletu równań różniczkowych opisujących stan układu dla każdej konfiguracji kluczy półprzewodnikowych.
	\item Opracowanie autorskiego silnika fizycznego symulatora opartego na numerycznej metodzie Eulera ze skrajnie krótkim krokiem całkowania ($\Delta t = 0{,}5~\mu\mathrm{s}$), z~wykorzystaniem wektoryzowanych operacji biblioteki NumPy~\cite{harris2020numpy}.
	\item Implementacja bezmodelowej logiki sterowania Bang-Bang (MF-BB) zdolnej radzić sobie z~dynamiką fazy nieminimalnej, generującej sygnał odniesienia prądu cewki bezpośrednio z~mierzonych sygnałów obwodowych~\cite{tatari2024mfbb}.
	\item Integracja z algorytmem PSO~\cite{kennedy1995pso, shi1998pso} w oparciu o bibliotekę naukową \texttt{PySwarms}~\cite{miranda2018pyswarms}, dedykowanym do automatycznej, dwutorowej optymalizacji parametrów sterownika MF-BB poprzez minimalizację ważonej funkcji kosztu, penalizującej uchyby napięcia oraz tętnienia prądu.
	\item Przeprowadzenie analizy porównawczej przebiegów uzyskanych z~opracowanego symulatora względem przebiegów referencyjnych wyeksportowanych z~programu PSIM, jak również serii testów badających odporność układu w~warunkach skrajnie zmniejszonej pojemności kondensatora wyjściowego.
\end{enumerate}

\section{Oczekiwane rezultaty}
\label{sec:rezultaty}
Zwieńczeniem tego projektu jest dostarczenie w pełni działającej aplikacji w~języku Python, zdolnej do precyzyjnego symulowania układu Boost z~rozdzielczością czasową rzędu pojedynczych mikrosekund. Publiczna dostępność kodu w systemie kontroli wersji Git~\cite{chacon2014progit} zapewnia pełną transparentność i~reprodukowalność uzyskanych wyników. Wiarygodność narzędzia została wykazana poprzez bezpośrednie numeryczne nałożenie sygnałów napięcia i~prądu na dane referencyjne z~oprogramowania PSIM. Część badawcza pracy wykazuje skuteczność metody PSO w~zautomatyzowanej redukcji tętnień napięcia wyjściowego oraz zapadów w~warunkach pracy z~małymi pojemnościami filtra, stanowiąc efektywną alternatywę dla manualnego strojenia kontrolerów.

\section{Założenia architektoniczne}
\label{sec:zalozenia}

Praca powstaje w~ramach kierunku Informatyka Stosowana, co determinuje szczególne wymagania stawiane wytwarzanemu kodowi. Naczelnym założeniem projektowym jest zastosowanie paradygmatu programowania obiektowego (ang. \emph{Object-Oriented Programming}, OOP) z~wyraźnym rozdziałem warstw odpowiedzialności:

\begin{itemize}
	\item \emph{warstwa fizyki obwodowej} -- klasy reprezentujące parametry, stany i~równania ruchu przekształtnika, niezależne od logiki sterowania;
	\item \emph{warstwa sterowania} -- klasy implementujące algorytmy regulacji (MF-BB), operujące wyłącznie na zdefiniowanych interfejsach pomiarowych;
	\item \emph{warstwa optymalizacji} -- moduły realizujące algorytm PSO we współpracy z biblioteką \texttt{PySwarms}, traktujące pętlę symulacyjną jako deterministyczną funkcję celu;
	\item \emph{warstwa eksperymentu i~wizualizacji} -- moduły uruchomieniowe oraz prezentacyjne, oddzielone od logiki rdzenia.
\end{itemize}

Taka separacja funkcjonalności umożliwia niezależną ewolucję komponentów, na przykład zastąpienie sterownika MF-BB regulatorem PI czy MPC, bez ingerencji w~pozostałe warstwy. Wykorzystanie systemu Git~\cite{chacon2014progit} do zarządzania kodem źródłowym pozwala na precyzyjne śledzenie ewolucji projektu oraz dokumentowanie kolejnych etapów eksperymentów.

\section{Struktura pracy}
\label{sec:struktura}

Struktura niniejszej pracy podzielona jest na sześć rozdziałów merytorycznych. Rozdział pierwszy (\nameref{ch:wstep}) wprowadza czytelnika w~tematykę pracy, formułuje cel oraz zakres prowadzonych badań. Rozdział drugi (\nameref{ch:analiza}) systematyzuje wiedzę o współczesnych metodach sterowania przekształtnikami, takimi jak techniki Bang-Bang~\cite{tatari2024mfbb} czy zaawansowane sterowanie predykcyjne~\cite{tatari2025splitcost}. Rozdział trzeci (\nameref{ch:realizacja}) stanowi rdzeń informatyczny pracy i~zawiera techniczny opis implementacji obiektowego modelu Bliźniaka Cyfrowego w języku Python. Rozdział czwarty (\nameref{ch:badania}) przedstawia metodykę badań symulacyjnych, empiryczną walidację modelu względem oprogramowania PSIM oraz stress-testy obciążeniowe i~pojemnościowe. Rozdział piąty (\nameref{ch:ai}) jest poświęcony integracji Bliźniaka Cyfrowego z~algorytmem Optymalizacji Rojem Cząstek (PSO), opisowi procedury dwutorowej oraz badaniu wrażliwości. Rozdział szósty (\nameref{ch:podsumowanie}) podsumowuje stopień realizacji celów inżynierskich, napotkane wyzwania oraz formułuje wnioski końcowe i~kierunki dalszych prac.
